\begin{pastexam}[label={px:ntust-110-ete-ds-07}]{臺科大電子所甲組 110 學年度 第 7 題}
（問答）

{[}10\%{]} Assume that a hash function (h(key)= key mod 8) is used for a chained hash table with
linked lists, where linked lists can be used to handle the collision due to hash function. If we
insert 10 items whose keys are 8, 9, 13, 17, 10, 15, 20, 16, 25 and 26 to an empty chained hash
table with linked lists, what are the execution results of the hash table?
\end{pastexam}

\begin{pastsoln}
表有 $8$ 個位置 $0$ 到 $7$，每個位置一條串列（鏈結法，資結書第~12.2~節）。 先算每個鍵的
$h(\mathit{key}) = \mathit{key} \bmod 8$：
\begin{center}
\begin{tabular}{l*{10}{c}}
\toprule
key & 8 & 9 & 13 & 17 & 10 & 15 & 20 & 16 & 25 & 26 \\
$h$ & 0 & 1 & 5 & 1 & 2 & 7 & 4 & 0 & 1 & 2 \\
\bottomrule
\end{tabular}
\end{center}
題目沒有規定新節點接在串列的頭還是尾。 下面照插入的先後接在\textbf{尾端}（串列裡由前到後就是插入
順序；斜線表示 \NIL）：
\begin{center}
\begin{tikzpicture}[>={Latex[length=4pt]},
  hd/.style={draw,minimum width=0.55cm,minimum height=0.5cm,inner sep=0pt},
  nd/.style={draw,minimum width=0.65cm,minimum height=0.5cm,inner sep=0pt,font=\small},
  lk/.style={draw,minimum width=0.4cm,minimum height=0.5cm,inner sep=0pt}]
  \node[font=\small,anchor=east] at (-0.35,-0.00) {[0]};
  \node[hd] (h0) at (0,-0.00) {};
  \fill (h0.center) circle (1.2pt);
  \node[nd] (n00) at (1.10,-0.00) {8};
  \node[lk,anchor=west] (l00) at (n00.east) {};
  \draw[->] (h0.center)--(n00.west);
  \node[nd] (n01) at (2.70,-0.00) {16};
  \node[lk,anchor=west] (l01) at (n01.east) {};
  \draw[->] (l00.center)--(n01.west);
  \fill (l00.center) circle (1.2pt);
  \draw (l01.south west)--(l01.north east);
  \node[font=\small,anchor=east] at (-0.35,-0.55) {[1]};
  \node[hd] (h1) at (0,-0.55) {};
  \fill (h1.center) circle (1.2pt);
  \node[nd] (n10) at (1.10,-0.55) {9};
  \node[lk,anchor=west] (l10) at (n10.east) {};
  \draw[->] (h1.center)--(n10.west);
  \node[nd] (n11) at (2.70,-0.55) {17};
  \node[lk,anchor=west] (l11) at (n11.east) {};
  \draw[->] (l10.center)--(n11.west);
  \fill (l10.center) circle (1.2pt);
  \node[nd] (n12) at (4.30,-0.55) {25};
  \node[lk,anchor=west] (l12) at (n12.east) {};
  \draw[->] (l11.center)--(n12.west);
  \fill (l11.center) circle (1.2pt);
  \draw (l12.south west)--(l12.north east);
  \node[font=\small,anchor=east] at (-0.35,-1.10) {[2]};
  \node[hd] (h2) at (0,-1.10) {};
  \fill (h2.center) circle (1.2pt);
  \node[nd] (n20) at (1.10,-1.10) {10};
  \node[lk,anchor=west] (l20) at (n20.east) {};
  \draw[->] (h2.center)--(n20.west);
  \node[nd] (n21) at (2.70,-1.10) {26};
  \node[lk,anchor=west] (l21) at (n21.east) {};
  \draw[->] (l20.center)--(n21.west);
  \fill (l20.center) circle (1.2pt);
  \draw (l21.south west)--(l21.north east);
  \node[font=\small,anchor=east] at (-0.35,-1.65) {[3]};
  \node[hd] (h3) at (0,-1.65) {};
  \draw (h3.south west)--(h3.north east);
  \node[font=\small,anchor=east] at (-0.35,-2.20) {[4]};
  \node[hd] (h4) at (0,-2.20) {};
  \fill (h4.center) circle (1.2pt);
  \node[nd] (n40) at (1.10,-2.20) {20};
  \node[lk,anchor=west] (l40) at (n40.east) {};
  \draw[->] (h4.center)--(n40.west);
  \draw (l40.south west)--(l40.north east);
  \node[font=\small,anchor=east] at (-0.35,-2.75) {[5]};
  \node[hd] (h5) at (0,-2.75) {};
  \fill (h5.center) circle (1.2pt);
  \node[nd] (n50) at (1.10,-2.75) {13};
  \node[lk,anchor=west] (l50) at (n50.east) {};
  \draw[->] (h5.center)--(n50.west);
  \draw (l50.south west)--(l50.north east);
  \node[font=\small,anchor=east] at (-0.35,-3.30) {[6]};
  \node[hd] (h6) at (0,-3.30) {};
  \draw (h6.south west)--(h6.north east);
  \node[font=\small,anchor=east] at (-0.35,-3.85) {[7]};
  \node[hd] (h7) at (0,-3.85) {};
  \fill (h7.center) circle (1.2pt);
  \node[nd] (n70) at (1.10,-3.85) {15};
  \node[lk,anchor=west] (l70) at (n70.east) {};
  \draw[->] (h7.center)--(n70.west);
  \draw (l70.south west)--(l70.north east);
\end{tikzpicture}
\end{center}
如果照資結書虛擬碼~12.1 把新節點插在串列\textbf{最前面}（插入 $O(1)$，不用走到尾端），每條串列的順序
反過來：位置 0 是 $16 \to 8$，位置 1 是 $25 \to 17 \to 9$，位置 2 是 $26 \to 10$，其餘相同。 兩種
都對，作答時寫明採用哪一種。

$10$ 個鍵放進 $8$ 個位置，載入因子 $\alpha = 10/8 = 1.25$。 最長的串列在位置 1，有 $3$ 個；位置 3、6
是空的。
\end{pastsoln}
